%======================================================================%
\section{The Mirror Climb}
\label{sec:climb}
%======================================================================%

\begin{equation}\label{eq:climb:unified}
  \Void\to\R\to\C\to\Quat\to\Oct\to\jalg\to\E{8}\to\E{8}\times\E{8}\to\E{6}\to\cdots\to\Void
\end{equation}

\begin{remark}[Climb sequence forced, not chosen]\label{rem:climb:forced}
The orbit \eqref{eq:climb:unified} is fixed by void, mirror, and capacity---not
by importing string-theory rungs. \emph{Functoriality}
(Definition~\ref{def:mirror}) restricts each step to modes
\textbf{M1}--\textbf{M4} of Definition~\ref{def:mirror:modes}. \emph{Hurwitz
terminal} at $n=4$ (Theorem~\ref{thm:rung-table}) forces the Jordan exit at
$n=5$ and anchors the associative subalgebra $\Quat\subset\Oct$ used below in
soldering (Definition~\ref{def:soldering}). On the stabilized Albert tower,
automorphism mirror closes at $\E{8}$ ($n=6$), product mirror is the unique
\textbf{M4} doubling on a simple Lie rung ($n=7$), and Cartan mirror is the
capacity-balanced collapse back to the $\jalg$ structure group ($n=8$,
Theorem~\ref{thm:e6}). \emph{Observer-slot requirement:} three independent
$\sigma$-anti-invariant classes in $H^1(\hat Q,V_{16})$ arise only after the
heterotic product quotient at $\Mirror^7(\Void)$ (Proposition~\ref{prop:climb:minimal-heterotic},
Section~\ref{sec:observer}); descent to $\E{6}$ supplies dynamics but does not
create the slot space. The above--below correspondence
(Definition~\ref{def:above-below}) ties each rung choice to a localized datum on
$\hat Q$.
\end{remark}

%----------------------------------------------------------------------%
\subsection{Division rungs}
\label{sec:climb:division}
%----------------------------------------------------------------------%

\begin{theorem}[Mirror rung table, $n\leq5$]\label{thm:rung-table}
\begin{center}
\renewcommand{\arraystretch}{1.15}
\begin{tabular}{@{}rcccc@{}}
\toprule
$n$ & $A_n$ & $d_n$ & Mode & \\
\midrule
$0$ & $\{0\}$ & $0$ & void & \\
$1$ & $\R$ & $1$ & discrete & \\
$2$ & $\C$ & $2$ & arithmetic & \\
$3$ & $\Quat$ & $4$ & arithmetic & \\
$4$ & $\Oct$ & $8$ & Hurwitz terminal & \\
$5$ & $\jalg$ & $27$ & Jordan exit & \\
\bottomrule
\end{tabular}
\end{center}
In particular $\Mirror^4(\Void)=\Oct$ and $\Mirror^5(\Void)\cong\jalg$.
\end{theorem}

\begin{proof}
Cayley--Dickson doubling until $\Oct$; Hurwitz uniqueness of normed division
algebras; sedenion zero divisors force Jordan exit to the unique formally real
$H_3(\Oct)$~\cite{Hurwitz1898,SpringerVeldkamp2000}.
\end{proof}

\begin{lemma}[Sedenion failure]\label{lem:sedenion}
$\Mirror_{\mathrm{arith}}(\Oct)$ is $16$-dimensional and not division.
\end{lemma}

%----------------------------------------------------------------------%
\subsection{Exceptional rungs}
\label{sec:climb:exceptional}
%----------------------------------------------------------------------%

On $\jalg$, the stabilizer tower is $G_2\subset F_4\subset\E{6}$. The
Freudenthal diagonal $\mathfrak{M}(\Oct,\Oct)=\mathfrak{e}_8$ yields:

\begin{theorem}[$\Mirror^6(\Void)=\E{8}$]\label{thm:e8}
Automorphism mirror on the octonion--Jordan pair gives $\E{8}$, $d_6=248$.
\end{theorem}

\begin{definition}[Lie product mirror]\label{def:prod-mirror}
On a stabilized exceptional rung with Lie group $G$, the \emph{product mirror}
is $\Mirror_{\mathrm{prod}}(G):=G\times G$ with involution
$\sigma:(g_1,g_2)\mapsto(g_2,g_1)$ exchanging factors---mode~\textbf{M4} at the
Lie level, functorially analogous to Boolean doubling at $n=0\to1$.
\end{definition}

\begin{theorem}[$\Mirror^7(\Void)=\E{8}\times\E{8}$]\label{thm:e8e8}
Product mirror $\Mirror_{\mathrm{prod}}(G)=G\times G$ gives $d_7=496$.
\end{theorem}

\begin{proposition}[Minimal heterotic rung for three observer slots]\label{prop:climb:minimal-heterotic}
$\Mirror^7(\Void)=\E{8}\times\E{8}$ is the \emph{minimal} mirror rung supporting
three observer slots: $\dim H^1(\hat Q,V_{16})^{\sigma=-1}=3$
(Theorem~\ref{thm:observer-fp}). At $n\leq6$ there is no product-mirror
$\mathbb{Z}_2$ quotient on the stabilized datum, hence no
$\Pi_{\sigma=-1}$ observer sector (Section~\ref{sec:observer:definition}). At
$n=8$ the heterotic doubling is already capacity-collapsed to $\E{6}$
(Theorem~\ref{thm:e6}), so the three $1$-classes are fixed on $\hat Q$ at
$n=7$ and only transported below by descent. Thus observer multiplicity is a
\emph{below}-level consequence of the \emph{above} product mirror at the
seventh rung.
\end{proposition}

\begin{proof}[Proof sketch]
For $n\leq6$, $\Mirror^n(\Void)$ is a single simple Lie datum or division
algebra; involution stabilizes within one chart, not across a heterotic pair.
For $n=7$, $\sigma$ acts on $Q$ and antiholomorphically on $V_{16}$; the
equivariant Lefschetz index on $H^1(Q,V_{16})$ counts three anti-invariant
contributions (Section~\ref{sec:observer:fixed-point}). For $n=8$,
$\Pi_{\mathrm{het}}:\E{8}\times\E{8}\to\E{6}$ preserves the slot basis on
$\hat Q$ while reducing the gauge group.
\end{proof}

\begin{theorem}[$\Mirror^8(\Void)=\E{6}$]\label{thm:e6}
Cartan mirror after capacity descent on the Albert graph gives $d_8=78$,
$H=\Spin(10)\times U(1)$, $\mathcal{M}=\E{6}/H$ (EIII), $\dim\mathcal{M}=32$.
\end{theorem}

\begin{proof}[Proof sketches]
\emph{$E_8$ (Steps 1--3).}
\emph{(i)} On $\jalg$, the stabilizer tower $G_2\subset F_4\subset\E{6}$ is
mirror-refined; the octonion measure before Jordan projection carries an
automorphism mirror $\Mirror_{\mathrm{aut}}$ coupling the two $\Oct$ sectors.
\emph{(ii)} Freudenthal--Tits identifies
$\mathfrak{M}(\Oct,\Oct)=\mathfrak{e}_8$ as the unique simply connected simple
Lie algebra closed under this coupling; $G_2\times F_4\subset\E{8}$ is maximal
among candidates containing the tower together with the mirrored octonion
sector~\cite{Freudenthal1954,Tits1966}.
\emph{(iii)} $d_6=248$; Stone chart dimension $\chi(S_6)=248$
(Definition~\ref{def:chart-dim}).

\emph{$E_8\times E_8$ (Steps 1--2).}
\emph{(i)} Apply $\Mirror_{\mathrm{prod}}$ (Definition~\ref{def:prod-mirror})
to the stabilized datum of Theorem~\ref{thm:e8}; functoriality of $\Mirror$
gives $\Mirror^7(\Void)=\Mirror(\Mirror^6(\Void))\cong\E{8}\times\E{8}$.
\emph{(ii)} $d_7=2\cdot248=496$; involutivity
(Proposition~\ref{prop:mirror:involution}) holds factorwise, with $\sigma$
implementing $\Mirror^2\cong\mathrm{id}$ up to diagonal gauge.

\emph{$E_6$ (Steps 1--3).}
\emph{(i)} On $\E{8}\times\E{8}$, the diagonal embedding
$\delta:\E{8}\hookrightarrow\E{8}\times\E{8}$ is the $\sigma$-stabilized
subgroup; $\mathrm{Bal}_{\mathcal{C}}$ on $\mathcal{G}_{\mathrm{Albert}}$ at
the $\jalg$ rung fixes the $1|10|16$ partition
(Theorem~\ref{thm:partition}).
\emph{(ii)} Among exceptional groups, only $\E{6}$ admits
$H=\Spin(10)\times U(1)$ with branching
$\mathbf{27}=\mathbf{1}\oplus\mathbf{10}\oplus\mathbf{16}$ and coset mirror
$\mathbf{16}\leftrightarrow\overline{\mathbf{16}}$; heterotic collapse is
indexed by the rank-$3$ diagonal of $\jalg$, not by the two $\E{8}$ factors
(Proposition~\ref{prop:climb:minimal-heterotic}).
\emph{(iii)} $d_8=\dim\E{6}=78$, $\mathcal{M}=\E{6}/H$ (EIII),
$\dim\mathcal{M}=32$.
\end{proof}

\begin{proposition}[Exceptional mirror steps]\label{prop:climb:exceptional}
At $n\geq6$ on the stabilized $\jalg$ rung, $\Mirror$ acts as:
\begin{enumerate}[label=(\roman*),leftmargin=*]
  \item $\E{8}$: automorphism mirror on the octonion--Jordan pair
    (Theorem~\ref{thm:e8});
  \item $\E{8}\times\E{8}$: Lie product mirror $\Mirror_{\mathrm{prod}}$
    (Theorem~\ref{thm:e8e8}, Definition~\ref{def:prod-mirror});
  \item $\E{6}$: Cartan mirror after capacity descent
    (Theorem~\ref{thm:e6});
  \item $\E{6}\to H=\Spin(10)\times U(1)$: coset stabilization on EIII;
  \item $\cdots\to\Void$: collapse after electroweak breaking.
\end{enumerate}
\end{proposition}

\begin{definition}[Chart dimension]\label{def:chart-dim}
At $n\geq6$, $\chi(S_n)=\dim A_n\in\{248,496,78\}$; profinite $|S_n|$ is
generally infinite. Stone refinement is chart synchronization, not cardinality
matching~\cite{Stone1936}.
\end{definition}